\documentclass[10pt]{scrartcl}
\usepackage[mdthm,colorsec,noauthor,noasy]{evan}
\usepackage{handout}
\DeclareMathOperator{\spn}{span}
\newcommand{\PP}{\mathcal P}
\title{Math 110 --- Week 2}
\subtitle{Solutions}
\author{}\date{}\ihead{}
\begin{document}
\maketitle

\begin{prob}\textbf{HW2 P1.}\par\idea{Sums can cancel, so look for the smallest cancellation.}
False. Take $n=1$, $v_1\ne0$ and $w_1=-v_1$. Each list is independent (one nonzero vector), but $v_1+w_1=0$ is a dependent list.
\end{prob}

\begin{prob}\textbf{HW2 P2.}\par\idea{(a) every axiom holds in each coordinate. (b) a finite-dimensional space cannot contain arbitrarily long independent lists. (c) finite combinations of the $e_j$ have only finitely many nonzero entries.}
\begin{steps}
  \ii (a) $\FF^\infty=\FF^{\mathbb N}$ with coordinatewise operations. Each vector space axiom is an equation between sequences, and it holds in every coordinate because it holds in $\FF$. Zero vector $(0,0,\dots)$; $-(a_j)=(-a_j)$.
  \ii (b) For every $m$, $e_1,\dots,e_m$ is independent: $\sum\lambda_je_j=(\lambda_1,\dots,\lambda_m,0,\dots)=0$ forces all $\lambda_j=0$. If $\FF^\infty$ were spanned by a list of length $n$, every independent list would have length $\le n$ (Prop.~10). Take $m=n+1$: contradiction.
  \ii (c) A finite combination $\lambda_1e_{j_1}+\dots+\lambda_ke_{j_k}$ is $0$ outside the positions $j_1,\dots,j_k$. The sequence $(1,1,1,\dots)$ has no zero entries, so it is not such a combination.
\end{steps}
\end{prob}

\begin{prob}\textbf{HW2 P3.}\par\idea{Start from a spanning list that lies in $U\cup W$ and thin it out; deleting vectors keeps the rest in $U\cup W$.}
\begin{steps}
  \ii $U$ and $W$ are finite-dimensional (Prop.~14). Let $u_1,\dots,u_k$ be a basis of $U$ and $w_1,\dots,w_l$ a basis of $W$.
  \ii The list $u_1,\dots,u_k,w_1,\dots,w_l$ spans $U+W=V$.
  \ii By Prop.~9 some of its vectors can be deleted to leave a basis of $V$. Every remaining vector is one of the $u_i$ or $w_j$, so it lies in $U\cup W$.
\end{steps}
\end{prob}

\begin{prob}\textbf{HW2 P4.}\par\idea{$\dim\FF^2=2$, so a proper nonzero subspace can only have dimension $1$, and two vectors form a basis exactly when they are independent.}
\begin{steps}
  \ii (a) Let $U$ be a subspace with $\{0\}\ne U\ne\FF^2$. By Prop.~14, $\dim U\le2$, and $\dim U\ne2$ since $U\ne\FF^2$; $\dim U\ne0$ since $U\ne\{0\}$. So $\dim U=1$: $U$ has a basis consisting of one vector $v\ne0$, and $U=\spn(v)$.
  \ii (b) ($\Rightarrow$) If $w=cv$, then $cv-w=0$ is a nontrivial relation, so $v,w$ is dependent; same if $v=cw$.
  \ii ($\Leftarrow$) Suppose neither is a multiple of the other, and $\lambda v+\mu w=0$. If $\mu\ne0$ then $w=-\frac\lambda\mu v$; if $\lambda\ne0$ then $v=-\frac\mu\lambda w$. Both are excluded, so $\lambda=\mu=0$. Two independent vectors in a $2$-dimensional space form a basis (Prop.~13).
\end{steps}
\end{prob}

\begin{prob}\textbf{HW2 P5.}\par\idea{The dimension formula turns ``$U+W$ is everything and the dimensions add up'' into ``the intersection is $0$''.}
$\dim(U\cap W)=\dim U+\dim W-\dim(U+W)=3+5-8=0$, so $U\cap W=\{0\}$ and the sum is direct.
\end{prob}

\begin{prob}\textbf{HW2 P6.}\par\idea{Both conditions are about the point $1$, so expand $p$ in powers of $x-1$: then $p(1)$ and $p'(1)$ are the first two coefficients.}
\begin{steps}
  \ii (a) The zero polynomial is in $U$. If $p,q\in U$: $(\lambda p+q)(1)=\lambda p(1)+q(1)=0$ and $(\lambda p+q)'(1)=\lambda p'(1)+q'(1)=0$.
  \ii (b) $1,(x-1),(x-1)^2,(x-1)^3,(x-1)^4$ is a basis of $\PP_4(\RR)$ (Ex.~18). Write $p=c_0+c_1(x-1)+c_2(x-1)^2+c_3(x-1)^3+c_4(x-1)^4$.
  \ii $p(1)=c_0$ and $p'(1)=c_1$. So $p\in U\iff c_0=c_1=0\iff p\in\spn\bigl((x-1)^2,(x-1)^3,(x-1)^4\bigr)$.
  \ii These three polynomials are independent (part of a basis). So they are a basis of $U$ and $\dim U=3$.
\end{steps}
\end{prob}

\begin{prob}\textbf{R1.}\par\idea{Four vectors in $\RR^4$: if they are independent they are a basis, and then both $U+W=\RR^4$ and $U\cap W=\{0\}$ follow.}
\begin{steps}
  \ii (a) Suppose $au_1+bu_2+cw_1+dw_2=0$. Coordinates: $a+d=0$, $a+b=0$, $b+c=0$, $c-d=0$. So $d=-a$, $b=-a$, $c=a$, and the last equation gives $a-(-a)=2a=0$. Hence $a=b=c=d=0$.
  \ii So $u_1,u_2,w_1,w_2$ is an independent list of $4=\dim\RR^4$ vectors, a basis (Prop.~13). Hence $U+W=\RR^4$, and $\dim U=\dim W=2$.
  \ii $\dim(U\cap W)=2+2-4=0$, so the sum is direct.
  \ii (b) Solve $au_1+bu_2+cw_1+dw_2=(2,1,3,0)$: $a+d=2$, $a+b=1$, $b+c=3$, $c-d=0$. From the last, $c=d$; then $b=3-d$, $a=2-d$, and $a+b=5-2d=1$ gives $d=2$. So $a=0$, $b=1$, $c=d=2$.
  \ii $u=u_2=(0,1,1,0)$ and $w=2w_1+2w_2=(2,0,2,0)$. Check: $u+w=(2,1,3,0)$.
\end{steps}
\end{prob}

\begin{prob}\textbf{R2.}\par\idea{Extend a basis of $U$ to a basis of $V$; the added vectors span the complement.}
\begin{steps}
  \ii Let $u_1,\dots,u_k$ be a basis of $U$ (Prop.~14) and extend it to a basis $u_1,\dots,u_k,w_1,\dots,w_l$ of $V$ (Prop.~11). Let $W=\spn(w_1,\dots,w_l)$.
  \ii $U+W$ contains a basis of $V$, so $U+W=V$.
  \ii If $x\in U\cap W$, then $x=\sum a_iu_i=\sum b_jw_j$, so $\sum a_iu_i-\sum b_jw_j=0$. The combined list is independent, so all $a_i,b_j$ are $0$ and $x=0$.
\end{steps}
\end{prob}

\begin{prob}\textbf{R3.}\par\idea{Write the condition as one linear equation in the coefficients and solve it with free parameters.}
\begin{steps}
  \ii For $p=a+bx+cx^2$: $p'(1)=b+2c$ and $p(2)=a+2b+4c$. The condition is $a+b+2c=0$, i.e.\ $a=-b-2c$.
  \ii So $p=b(x-1)+c(x^2-2)$ with $b,c$ free: $V=\spn(x-1,\ x^2-2)$. Check: for $x-1$, $p'(1)=1=p(2)$; for $x^2-2$, $p'(1)=2=p(2)$.
  \ii They are independent (different degrees). So $x-1,\ x^2-2$ is a basis and $\dim V=2$.
\end{steps}
\end{prob}

\end{document}
